\documentclass[12pt,a4paper]{article}

% --- Packages ---
\usepackage[utf8]{inputenc}
\usepackage[T1]{fontenc}
\usepackage[margin=1in]{geometry}
\usepackage{setspace}
\usepackage{hyperref}
\usepackage{amsmath, amssymb, amsfonts}
\usepackage{graphicx}
\usepackage{booktabs}
\usepackage{float}
\usepackage{caption}
\usepackage{subcaption}
\usepackage{cite}
\bibliographystyle{unsrt}

% --- Formatting ---
\setstretch{1.2}
\hypersetup{colorlinks=true, linkcolor=blue, citecolor=blue, urlcolor=blue}

% --- Title Information ---
\title{\Large \textbf{Hierarchical Edge-Cloud Ensemble Learning for Safety-Critical Poultry Health Monitoring: A Recall-First Approach}}
\author{\textbf{Heng SUN} \\ 
        \textit{Bachelor of Data Science and AI Engineering} \\
        \textit{Cambodia University of Technology and Science (CamTech)}}
\date{March 15, 2026}

\begin{document}

\maketitle

\begin{abstract}
Automated disease detection in poultry farming is often hindered by a binary trade-off: the latency of high-precision cloud models versus the hardware constraints of on-farm edge devices. This paper proposes a hierarchical "Gatekeeper" architecture designed for the Tokkatot Smart Farming ecosystem. We implement a two-stage pipeline: a lightweight YOLOv8 screening model deployed on Raspberry Pi for 24/7 manure conveyor monitoring, and a high-fidelity Cloud Ensemble composed of EfficientNetB0 and DenseNet121. To address the asymmetric cost of diagnostic errors in agriculture, we introduce a \textit{Recall-Focused Loss} ($\lambda=5$) and a \textit{Safety-First Voting Gate}. Our results demonstrate a 99.1\% accuracy on classified samples and a near-perfect disease recall (0.998), while successfully isolating 5.01\% of ambiguous cases for manual review. This research provides a scalable framework for deploying mission-critical deep learning in resource-constrained agricultural environments.
\end{abstract}

\newpage

\section{Introduction}
The global poultry industry faces a persistent threat from highly transmissible diseases such as Coccidiosis, Newcastle Disease, and Salmonella. In Cambodia, where medium and small-scale poultry farming forms the backbone of rural livelihoods, these outbreaks are catastrophic. A single undetected infection can lead to the loss of an entire flock within 72 hours, resulting in devastating financial losses for farmers who often operate with thin margins and limited access to professional veterinary services.

\subsection{Tokkatot: From Student Vision to Smart Solution}
Tokkatot originated as a mission-driven initiative by a team of dedicated students at the Cambodia University of Technology and Science (CamTech). Witnessing the recurring cycle of disease-driven bankruptcy among local farmers, the team sought to bridge the gap between traditional agricultural practices and modern technology. The result is the \textbf{Tokkatot Smart Chicken Farming Solution}---an automated, IoT-based ecosystem designed specifically for the Cambodian context.

The broader Tokkatot platform aims to modernize traditional chicken raising by automating environmental control, feeding, and health monitoring. By reducing the heavy reliance on manual labor and providing real-time data insights, Tokkatot empowers small and medium farmers to transition from reactive crisis management to proactive, data-driven farming.

\subsection{Research Focus and Scope}
While the Tokkatot ecosystem encompasses various IoT modules, this paper focuses solely on the design and implementation of its most critical feature: the \textbf{Automated Chicken Disease Detection System}. We detail how this system was built from the ground up using the Poultry Diseases dataset \cite{natesh2024poultry}---from data collection in local environments to the development of a hierarchical Edge-Cloud deep learning architecture. Furthermore, we discuss the practical implementation strategies required to deploy this high-precision diagnostic tool on real-world Cambodian farms, where hardware cost and internet stability are primary constraints.

\section{System Architecture: The Hierarchical Gatekeeper}
The primary innovation of our system is the split between Edge Screening and Cloud Diagnosis. This hierarchical flow, illustrated in Figure \ref{fig:overview}, ensures that the system remains responsive and cost-effective.

\begin{figure}[H]
    \centering
    \includegraphics[width=0.9\textwidth]{figures/architecture_overview.png}
    \caption{The Tokkatot Hierarchical Architecture. The process begins with 24/7 YOLOv8 monitoring on the Edge. Only suspicious samples trigger the high-bandwidth upload to the Cloud Ensemble for multi-class diagnostic voting.}
    \label{fig:overview}
\end{figure}

\subsection{Edge Layer: YOLOv8 Gatekeeper}
The Edge layer utilizes a custom-trained, INT8-quantized YOLOv8 model \cite{jocher2023ultralytics}. Unlike classification-only models, YOLOv8 allows the system to localize specific fecal samples on a moving manure conveyor belt. The model is trained to perform binary screening: if a sample is classified as "Healthy" with high confidence, no further action is taken. This reduces cloud compute costs by over 90\% in a typical healthy flock.

\subsection{Cloud Layer: Safety-First Ensemble}
Once a sample is flagged as "Unhealthy" or "Suspicious" by the Edge, a high-resolution image is transmitted to the Tokkatot Cloud. Here, we employ a dual-model ensemble. We chose \textbf{EfficientNetB0} \cite{tan2019efficientnet} for its superior compound-scaling properties and \textbf{DenseNet121} \cite{huang2017densely} for its dense connectivity, which excels at capturing the mucosal textures of diseased feces. 

\section{Methodology and Training Strategy}

\subsection{Asymmetric Risk Handling: Recall-Focused Loss}
Standard Cross-Entropy loss assumes that all misclassifications carry equal weight. However, in poultry farming, the cost of a False Negative (FN) is the value of the entire flock, while the cost of a False Positive (FP) is merely a few minutes of human inspection. To reflect this, we modified the training objective to include a penalty multiplier $\lambda$:

\begin{equation}
\mathcal{L}(x, y) = 
\begin{cases} 
\lambda \cdot \text{CE}(x, y), & \text{if } y \in \{\text{Diseases}\} \text{ and } \hat{y} = \text{Healthy} \\
\text{CE}(x, y), & \text{otherwise}
\end{cases}
\end{equation}

By setting $\lambda=5$, we forced the backbones to be "hyper-sensitive" to disease features during the fine-tuning phase.

\subsection{Ensemble Consensus Logic}
The final diagnostic output is determined by a \textbf{Safety-First Voting Gate} (see Figure \ref{fig:logic}). This gate is more conservative than a standard average or majority vote. It requires:
\begin{enumerate}
    \item \textbf{High Confidence:} Both models must return a softmax probability $P > 0.85$.
    \item \textbf{Agreement:} Both models must predict the same class.
\end{enumerate}
If either condition fails, the sample is automatically routed to "Uncertain/Isolate" mode.

\begin{figure}[H]
    \centering
    \includegraphics[width=0.7\textwidth]{figures/ensemble_logic_safety_first.png}
    \caption{The internal logic of the Voting Gate. The system defaults to the safest action (Isolation) in any case of model disagreement or low confidence.}
    \label{fig:logic}
\end{figure}

\section{Results and Evaluation}

\subsection{Statistical Performance}
The ensemble was evaluated on a test set of 70,677 samples. As shown in Figure \ref{fig:performance_comparison}, the Ensemble consistently outperformed the individual backbones across all metrics. By combining the different inductive biases of EfficientNet and DenseNet, we achieved a more stable decision boundary.

\begin{figure}[H]
    \centering
    \includegraphics[width=0.8\textwidth]{figures/model_performance_comparison.png}
    \caption{Performance comparison between the standalone backbones and the Ensemble. Note the jump in F1-score and Accuracy provided by the voting consensus.}
    \label{fig:performance_comparison}
\end{figure}

\subsection{Safety and Isolation Rate}
A critical metric for farm operations is the \textbf{Isolation Rate}. We found that 5.01\% of samples triggered the safety gate's "Uncertain" mode (Figure \ref{fig:isolation}). Analysis of these samples revealed they were often low-quality images or samples with overlapping fecal matter, where a single model's prediction would have been statistically risky.

\begin{figure}[H]
    \centering
    \includegraphics[width=0.7\textwidth]{figures/isolation_statistics.png}
    \caption{The "5.01\% Rule." By isolating a small fraction of ambiguous cases, we ensure the remaining 94.99\% of automated classifications are extremely high-fidelity.}
    \label{fig:isolation}
\end{figure}

\subsection{Diagnostic Precision: Confusion Matrix Analysis}
The confusion matrices in Figure \ref{fig:confusion_matrices} highlight the system's ability to distinguish between visually similar diseases. While individual models occasionally confused Salmonella with Coccidiosis due to similar mucosal textures, the Ensemble's agreement rule successfully filtered these errors into the "Uncertain" category rather than committing to a wrong diagnosis.

\begin{figure}[H]
    \centering
    \begin{subfigure}[b]{0.31\textwidth}
        \includegraphics[width=\textwidth]{figures/confusion_matrix_efficientnet.png}
        \caption{EfficientNetB0}
    \end{subfigure}
    \hfill
    \begin{subfigure}[b]{0.31\textwidth}
        \includegraphics[width=\textwidth]{figures/confusion_matrix_densenet.png}
        \caption{DenseNet121}
    \end{subfigure}
    \hfill
    \begin{subfigure}[b]{0.31\textwidth}
        \includegraphics[width=\textwidth]{figures/confusion_matrix_ensemble.png}
        \caption{Ensemble}
    \end{subfigure}
    \caption{Comparison of confusion matrices. The Ensemble (right) shows a significant reduction in off-diagonal errors, particularly in the Salmonella/Coccidiosis overlap.}
    \label{fig:confusion_matrices}
\end{figure}

\section{Discussion: Engineering for the Real World}
The primary finding of this research is that \textbf{Recall-First} design is mandatory for agricultural AI. In the field, lighting conditions fluctuate from morning sun to evening fluorescent bulbs, and dust often accumulates on camera lenses. Our aggressive augmentation strategy and the ensemble's "I don't know" capability (Isolation mode) directly address these environmental noises.

Compared to existing state-of-the-art models that attempt to run heavy transformers on edge hardware, our hierarchical YOLO-Ensemble approach maintains a high frame rate (22+ FPS) on the Raspberry Pi while providing "Cloud-Grade" diagnostic precision when it matters most.

\section{Conclusion}
The Tokkatot AI module demonstrates a successful path for deploying safety-critical deep learning in the poultry sector. By employing a hierarchical Edge-Cloud pipeline, we successfully navigated the constraints of low-power hardware and high-precision requirements. The combination of \textit{Recall-Focused Loss} and \textit{Ensemble Safety Voting} ensures that the system serves as a reliable decision-support tool for farmers, significantly reducing the risk of catastrophic flock loss. Future work will involve expanding the dataset to include varying stages of disease progression and testing the system's robustness against different chicken breeds in Cambodia.

\section{References}
\bibliography{reference}

\end{document}
